As seen from an inertial reference frame outside Earth, the observer must
describe a circle parallel to the ecliptic, with constant speed, in a period
of one sidereal day. Since the ecliptic poles are in zero altitude, the circle
will be a great circle.

But the resulting movement should be the actual movement of the observer (in
relation of Earth's surface) plus the movement of Earth (in relation to the
inertial reference frame). After some time $t$, the observer must describe an
arc of length $\phi = 2\pi/T$ % sic: presumably 2*pi*t/T is meant
on the great circle; at the same time, the Earth will rotate by this same
amount.

In other words, the displacement of the observer is the composition of a
rotation with an angle $\varphi$ around the ecliptic axis and a rotation with
an angle $-\varphi$ around the earth's rotation axis. But this is exactly the
same construction which describes an analemma of a planet with a circular
orbit, inclination of $\epsilon = 23.44^\circ$ between the ecliptic and
equator!

So, the important characteristics of the drawing curve are:
\begin{description}
\item[a.] It is closed
\item[b.] Maximum and minimum latitudes are $\theta = \pm 23.44^\circ$
\item[c.] It is symmetric around the Equator
\item[d.] It is symmetric around a meridian
\item[e.] It has the shape of an eight (8) \hfill(1 pt each item)
\end{description}

For the calculation: as the observer starts on the southern hemisphere, he
will cross the Equator for the first time in the direction south--north. Its
velocity relative to an inertial referential will have intensity
$v = 2\pi R/T$ and azimuth (the angle with the north-south direction)
$\alpha = 90^\circ - \epsilon$. The velocity of the Earth's surface will have
the same intensity and azimuth $\alpha = 90^\circ$. Subtracting the vectors,
we'll have a resultant vector with modulus
\[
  v' = v\sqrt{2(1 - \cos\epsilon)}
\]
\hfill(3 pt velocity)

And azimuth
\[
  \alpha' = -\frac{\epsilon}{2} = -11.72^\circ
\]
\hfill(2 pt direction)
